\documentclass[10pt]{report}
\setcounter{tocdepth}{3}
\setcounter{secnumdepth}{3}
\usepackage[french]{babel}
\usepackage[a4paper, left=30mm,right=20mm,top=25mm,bottom=25mm]{geometry}
\usepackage{graphicx} % Required for inserting images
\usepackage[T1]{fontenc}

\title{Rapport CO (TCPP2)}
\date{November 2024}
\setlength{\parskip}{1ex plus 0.5ex minus 0.2ex}
\newcommand{\hsp}{\hspace{20pt}}
\newcommand{\HRule}{\rule{\linewidth}{0.5mm}}
\renewcommand\thesection{\arabic{section}}

\begin{document}

\begin{titlepage}
  %\begin{sffamily}
  \begin{center}
    \includegraphics[height=25mm]{gfx/OIP.jpg}~\\[3cm]

    \textsc{\LARGE UFR Sciences et technique}\\
    \textsc{\Large Université de Marine Le pen }\\
    \textsc{\Large \textbf{Projet}}\\
    \textsc{\Large L3 Informatique}\\[2.5cm]

    \HRule \\[0.4cm]
    { \huge \bfseries Jeu de plate-formes avec génération procédurale\\[0.4cm] } % C'est le titre donné par jube, j'y suis pour rien

    \HRule \\[0.6cm]\renewcommand\thesection{\arabic{section}}
    \textsc{\Large \textbf{Rapport de projet}}\\[1.5cm]
    \begin{minipage}{\textwidth}
      \begin{center}\LARGE
        Kilian \textsc{Jelic} \\
        Laura \textsc{Jacqueson} \\
        Théo \textsc{Pariney} \\
      \end{center}
    \end{minipage}
    \vfill
    % Bottom of the page
    {\large \ avril 2025}

  \end{center}
  %\end{sffamily}
\end{titlepage}

\normalsize
\pagenumbering{arabic}
\tableofcontents
\pagebreak
\listoffigures
\pagebreak

\section{Introduction}
%Le sujet du projet
%Pourquoi ce projet
%Qui a réalisé le projet
%Présentation du plan

\section{Présentation du jeu}
\subsection{Présentation générale}
\subsection{Les objectifs du joueur}
\subsubsection{Les écrous}
\subsubsection{La sortie}
\subsection{Les scènes}

\section{Les blocks et entités du monde}
\subsection{La définition dynamique des blocks}
% Fichier XML et récupération dans une classe dédiée
\subsection{Le stockage des blocks}
% Récupération dans le blockmanager utilisé par le monde

\section{La génération procédurale}
% Je sais pas trop comment ça se passe, faudrait le remplir
% Faudra penser à préciser qu'on a commencé par faire un monde "test" avant d'arriver à notre générateur actuel
% Et aussi à donner les étapes/difficultés avant d'en arriver à un truc qui marche correctement

\section{Le personnage}
\subsection{Les contrôles de base}
\subsection{Le saut}
\subsubsection{Le saut modulable}
%ça fait partie des choix qu'on a fait puis abandonné, c'est important à mentionner
\subsubsection{Le double saut}
\subsection{Le dash}
%Faudra aussi mentionner qu'on a eu du mal à trouver des contrôles convenables
\subsection{Le score et les vies}

\section{La physique}
\subsection{La gestion des collisions}
\subsubsection{Les propriétés des blocks}
\subsection{Calcul de la résistance de l'air}

\section{Notre organisation}
\subsection{La définition des tâches}
\subsubsection{La répartition du travail}
\subsection{Les outils utilisés}
% ça serait important de mentionner qu'on a mis un peu de temps à avoir un trello
% je pense aussi que l'ordre dans lequel on a mis les outils en place peut être interessant à rajouter
\subsubsection{Discord}
\subsubsection{Google Docs} %Pour les idées à la base
\subsubsection{Trello}
\subsubsection{Git et Github}
\subsubsection{IDE}
% Cette section là elle est marrante parse que personne utilise le même
\newpage

\section{Conclusion}
% On a quand même une conclusion ici avant les autres points
\subsection{Points d'amélioration}
\subsubsection{Ce qu'on aurait voulu ajouter}
%c.f le gdoc
\subsubsection{Wave Function Collapse}
\subsubsection{Pathfinding amélioré}
\subsubsection{La physique} % ;-;


\end{document}
